% Generated by scripts/assemble.py; corrections belong in source/corrections.json.
\examhead{148}{湖南师范大学 2025 年数学分析试卷(回忆版)}{140}
% p141-b23; source PDF page 141
\par\Needspace{5\baselineskip}
\blockanchor{p141-b23}
\majorquestion{一}填空题.\par
% p141-b24; source PDF page 141
\blockanchor{p141-b24}
\question{1}$\lim_{n\to \infty}\frac{\sum_{k = 1}^{n}\sqrt[k]{k^2 + 2^k}}{n} =$\par
% p142-b00; source PDF page 142
\blockanchor{p142-b00}
\question{2}$\iint_{[0,1]\times [0,1]}\left|xy - \frac{1}{4}\right|\mathrm{d}x\mathrm{d}y =$\par
% p142-b01; source PDF page 142
\blockanchor{p142-b01}
\question{3}已知 $f(x)$ 为连续函数, 且 $\lim_{x \to 0} \frac{\ln(1 + 2x^2) - xf(x)}{x^2} = 0$, 求 $\lim_{x \to 0} \frac{f(x)}{x} =$ \_\_\_\_.\par
% p142-b02; source PDF page 142
\blockanchor{p142-b02}
\question{4}设  $x_{n}=(1+\alpha)(1+\alpha^{2})(1+\alpha^{4})\cdots(1+\alpha^{2^{n}})$ ,  $|\alpha|<1$ , 求  $\lim_{n\to\infty}x_{n}=$  \_\_\_\_.\par
% p142-b03; source PDF page 142
\blockanchor{p142-b03}
\question{5}$\sup_{0 < s < 1}\int_0^1 |\ln |s - t||\mathrm{d}t =$\par
% p142-b04; source PDF page 142
\blockanchor{p142-b04}
\question{6}设 $f(x) = \begin{cases} x, & -1 \leq x \leq 2; \\ 0, & \text{其他}. \end{cases}$ 则 $\iint_{\mathbb{R}^2} f(x) f(x^2 - y) \, \mathrm{d}x \, \mathrm{d}y =$\par
% p142-b05; source PDF page 142
\blockanchor{p142-b05}
\question{7}设 $L: |x| + |y| = 1$, 取逆时针方向, 求 $\oint_{L} \frac{-y - x^{2}}{x^{2} + y^{2} + 2|xy|} \mathrm{d}x + \frac{x + y^{2}}{x^{2} + y^{2} + 2|xy|} \mathrm{d}y =$ \_\_\_\_.\par
% p142-b06; source PDF page 142
\blockanchor{p142-b06}
\question{8}$\lim_{n\to\infty}\frac{\cos\frac{\pi}{4n}+\cos\frac{3\pi}{4n}+\cdots+\cos\frac{(2n-1)\pi}{4n}}{\sqrt{n^{2}+n+1}}=$  \_\_\_\_.\par
% p142-b07; source PDF page 142
\blockanchor{p142-b07}
\question{9}级数  $\sum_{n=1}^{\infty}\frac{(-1)^{n}n+x\sin nx}{n^{2}+x^{2}}$  的收敛域为 \_\_\_\_.\par
% p142-b08; source PDF page 142
\blockanchor{p142-b08}
\question{10}求 $\iint_{S} |z| \mathrm{~d} S =$ \_\_\_\_，其中 $S$ 为柱体 $x^{2} + y^{2} \leq ax$ 与球体 $x^{2} + y^{2} + z^{2} \leq a^{2}$ 相交的曲面，$a > 0$.\par
% p142-b09; source PDF page 142
\par\Needspace{5\baselineskip}
\blockanchor{p142-b09}
\majorquestion{二}计算题.\par
% p142-b10; source PDF page 142
\blockanchor{p142-b10}
\question{1}计算广义积分 $\int_0^{+\infty} e^{-2x} |\sin x| \mathrm{d}x$.\par
% p142-b11; source PDF page 142
\blockanchor{p142-b11}
\question{2}设 $f(x) = \frac{1 - x + x^2}{1 + x + x^2}$，求 $f^{(2024)}(0)$.\par
% p142-b12; source PDF page 142
\par\Needspace{6\baselineskip}
\blockanchor{p142-b12}
\question{3}用 $C[-1,1]$ 表示闭区间 $[-1,1]$ 上连续函数的合集, 对任意的 $f(x) \in C[-1,1]$, 记\par
\nopagebreak[4]
% p142-b13; source PDF page 142
\blockanchor{p142-b13}
\displayformula{A (f (x)) = \int_ {- 1} ^ {0} f (x) \mathrm{d} x - \int_ {0} ^ {1} f (x) \mathrm{d} x.}
% p142-b14; source PDF page 142
\blockanchor{p142-b14}
求 $\sup \{A(f(x)) | \max_{-1 < x < 1} |f(x)| \leq 1, f(x) \in C[-1, 1]\}$.\par
% p142-b15; source PDF page 142
\par\Needspace{5\baselineskip}
\blockanchor{p142-b15}
\majorquestion{三}证明题.\par
% p142-b16; source PDF page 142
\par\Needspace{6\baselineskip}
\blockanchor{p142-b16}
\question{1}设 $f(x)$ 在任意的 $[0, b] (b > 0)$ 上的黎曼可积，且 $\lim_{x \to +\infty} f(x) = \alpha$。证明：\par
\nopagebreak[4]
% p142-b17; source PDF page 142
\blockanchor{p142-b17}
\displayformula{\lim _ {t \rightarrow 0 ^ {+}} t \int_ {0} ^ {+ \infty} e ^ {- t x} f (x) \mathrm{d} x = \alpha .}
% p142-b18; source PDF page 142
\blockanchor{p142-b18}
\question{2}设 $f(x)$ 在 $[a, b]$ 上有界, 在 $[a, b]$ 上只有 $\{x_n\}_{n=1}^{\infty}$ 为不连续点, 且 $\{x_n\}$ 收敛. 证明: $f(x)$ 在 $[a, b]$ 上可积.\par
% p142-b19; source PDF page 142
\blockanchor{p142-b19}
\question{3}证明 $\sin \frac{1}{x}$ 在 $(c,1)(0 < c < 1)$ 上一致连续，在 $(0,1)$ 上非一致连续.\par
% p142-b20; source PDF page 142
\blockanchor{p142-b20}
\question{4}描述区间套定理, 并用区间套定理证明致密性定理.\par
